% One-time imported draft. Editable native source; no Markdown regeneration.
\section{Step 15 Design a custom
  glyph}\label{005-seven-segment-counter-s15-step-15-design-a-custom-glyph}

Bit 7 controls the decimal point independently of the seven bars.
Turning it on for the A mask changes 77 to f7. In a source byte
initializer, the signed value -9 expresses that pattern without a
truncation warning.

\subsection{Write or
  inspect}\label{005-seven-segment-counter-s15-write-or-inspect}

In a separate copy, change the A table entry to -9 and predict its
display. Keep the submitted baseline table unless your assignment
intentionally uses the custom design.

\subsection{Expected
  observation}\label{005-seven-segment-counter-s15-expected-observation}

The A bars remain and the decimal point is added. The bus value is the
unsigned byte 247, regardless of the signed initializer used to store
its bits.

\subsection{Check your
  understanding}\label{005-seven-segment-counter-s15-check-your-understanding}

What does setting bit 7 change?

\begin{itemize}
  \tightlist
  \item
        \textbf{A} The number of loop iterations.
  \item
        \textbf{B} The decimal point on that display.
  \item
        \textbf{C} The address of the left display.
\end{itemize}
